\section{实验设计}

\subsection{数据通路}\label{sub:datapath}

\subsubsection{功能描述}
数据通路（Datapath）负责完成指令执行过程中所有数据的流动与运算，支持R型指令（add、sub、and、or、slt）、立即数指令（addi）、访存指令（lw、sw）、分支指令（beq）和跳转指令（j）共十条指令的执行。

\subsubsection{接口定义}

\begin{table}[htp]
\caption{Datapath接口定义}\label{tab:datapath}
\begin{center}
\begin{tabular}{|l|l|l|p{6cm}|}
\hline
\textbf{信号名} & \textbf{方向} & \textbf{位宽} & \textbf{功能描述}\\ \hline \hline
clka         & Input  & 1-bit  & 时钟信号 \\ \hline
rst          & Input  & 1-bit  & 同步复位信号，高电平有效 \\ \hline
jump         & Input  & 1-bit  & 跳转控制信号（j指令） \\ \hline
branch       & Input  & 1-bit  & 分支控制信号（beq指令） \\ \hline
alusrc       & Input  & 1-bit  & ALU第二操作数选择：0=寄存器，1=立即数 \\ \hline
memtoreg     & Input  & 1-bit  & 写回数据选择：0=ALU结果，1=存储器读出数据 \\ \hline
regwrite     & Input  & 1-bit  & 寄存器堆写使能 \\ \hline
regdst       & Input  & 1-bit  & 目的寄存器选择：0=rt，1=rd \\ \hline
alucontrol   & Input  & 3-bit  & ALU操作控制信号 \\ \hline
pc           & Output & 32-bit & 当前程序计数器值 \\ \hline
instr        & Input  & 32-bit & 从指令存储器读出的指令 \\ \hline
aluout       & Output & 32-bit & ALU运算结果（兼作数据存储器地址） \\ \hline
writedata    & Output & 32-bit & 写入数据存储器的数据（来自寄存器rt） \\ \hline
readdata     & Input  & 32-bit & 从数据存储器读出的数据 \\ \hline
display\_data & Output & 32-bit & \$s0寄存器（reg[16]）值，供七段数码管显示 \\ \hline
\end{tabular}
\end{center}
\end{table}

\subsubsection{逻辑控制}
数据通路依照指令执行的五个阶段依次连接各子模块：
\begin{enumerate}[(1)]
    \item \textbf{取指（IF）}：PC寄存器输出当前地址，指令存储器根据地址读出指令，同时加法器计算PC+4。
    \item \textbf{译码（ID）}：指令字段分别送往寄存器堆和控制器，寄存器堆读出rs、rt对应的数据RD1、RD2；符号扩展模块对立即数字段进行16位到32位的符号扩展。
    \item \textbf{执行（EX）}：ALU根据alucontrol对SrcA、SrcB进行运算，其中SrcB由ALUSrc信号在RD2与符号扩展立即数之间选择；对于beq指令，将符号扩展立即数左移2位后与PC+4相加得到分支目标地址PCBranch；对于j指令，将instr[25:0]左移2位后拼接PC+4[31:28]得到PCJump。
    \item \textbf{访存（MEM）}：lw指令从数据存储器读取数据，sw指令将寄存器数据写入存储器；beq指令通过ALU zero信号与branch信号做AND运算得到PCSrc，控制下一PC。
    \item \textbf{写回（WB）}：通过MemtoReg信号在ALU结果与存储器读出数据之间选择，将结果写入目的寄存器；目的寄存器地址由RegDst信号在rt与rd之间选择。
\end{enumerate}

\subsection{控制器}\label{sub:controller}

\subsubsection{功能描述}
控制器接收32位指令，提取opcode和funct字段，经主译码器和ALU译码器生成全部控制信号，驱动数据通路完成相应指令的执行。

\subsubsection{接口定义}

\begin{table}[htp]
\caption{Controller接口定义}\label{tab:controller}
\begin{center}
\begin{tabular}{|l|l|l|p{6cm}|}
\hline
\textbf{信号名} & \textbf{方向} & \textbf{位宽} & \textbf{功能描述}\\ \hline \hline
inst       & Input  & 32-bit & 输入指令字 \\ \hline
memtoreg   & Output & 1-bit  & 写回数据选择控制 \\ \hline
memwrite   & Output & 1-bit  & 数据存储器写使能 \\ \hline
branch     & Output & 1-bit  & 分支控制信号 \\ \hline
alusrc     & Output & 1-bit  & ALU源操作数选择 \\ \hline
regdst     & Output & 1-bit  & 目的寄存器选择 \\ \hline
regwrite   & Output & 1-bit  & 寄存器堆写使能 \\ \hline
jump       & Output & 1-bit  & 跳转控制信号 \\ \hline
alucontrol & Output & 3-bit  & ALU操作控制信号 \\ \hline
\end{tabular}
\end{center}
\end{table}

\subsubsection{逻辑控制}
控制器由主译码器（maindec）和ALU译码器（aludec）两级组成。主译码器根据opcode（instr[31:26]）生成除alucontrol以外的全部控制信号及2位中间信号aluop；ALU译码器根据aluop与funct字段（instr[5:0]）生成alucontrol[2:0]。各指令对应的控制信号如表\ref{tab:truth}所示。

\begin{table}[htp]
\caption{控制信号真值表}\label{tab:truth}
\begin{center}
\begin{tabular}{|l|l|l|l|l|l|l|l|l|}
\hline
\textbf{指令} & \textbf{regwrite} & \textbf{regdst} & \textbf{alusrc} & \textbf{branch} & \textbf{memwrite} & \textbf{memtoreg} & \textbf{jump} & \textbf{aluop} \\ \hline \hline
R型   & 1 & 1 & 0 & 0 & 0 & 0 & 0 & 10 \\ \hline
lw    & 1 & 0 & 1 & 0 & 0 & 1 & 0 & 00 \\ \hline
sw    & 0 & 0 & 1 & 0 & 1 & 0 & 0 & 00 \\ \hline
beq   & 0 & 0 & 0 & 1 & 0 & 0 & 0 & 01 \\ \hline
addi  & 1 & 0 & 1 & 0 & 0 & 0 & 0 & 00 \\ \hline
j     & 0 & 0 & 0 & 0 & 0 & 0 & 1 & 00 \\ \hline
\end{tabular}
\end{center}
\end{table}
